\chapter{MÓDULO DE EVALUACIÓN}

Este módulo es donde se rinde el Desafío. El estudiante entra con su código personal, resuelve las preguntas de su categoría en el tiempo fijado y entrega; el sistema corrige cada respuesta en el servidor, calcula los puntajes y el ranking de cada categoría y publica los resultados en la fecha prevista. En las aulas sin conectividad, la prueba se rinde en papel y el maestro carga después las respuestas, que se corrigen igual.

\section{PLANIFICACIÓN DEL MÓDULO}

Esta sección fija el objetivo del módulo, nombra sus historias de usuario y reparte el trabajo en dos iteraciones (Tabla~\ref{tab:iteraciones_evaluacion}).

\subsection{Objetivo del módulo}

Soportar la resolución de los desafíos en entorno digital y la generación de formatos imprimibles para aulas sin conectividad, junto con la corrección de las respuestas y la generación de resultados.

\subsection{Historias de usuario}

Las historias responden al estudiante, que rinde; al administrador, que publica los resultados; y al maestro, que sigue a su grupo durante la prueba, imprime la prueba para las aulas sin conexión, carga las hojas de respuestas y revisa los resultados. Las tarjetas completas están en el Anexo A.

{\estilotabla
\begin{longtable}{l >{\raggedright\arraybackslash}p{6.9cm} >{\raggedright\arraybackslash}p{2.4cm} c c}
\arrayrulecolor{huazul}\specialrule{1.4pt}{0pt}{0pt}
\rowcolor{hufondo}\textbf{ID} & \textbf{Historia} & \textbf{Usuario} & \textbf{Prioridad} & \textbf{Esfuerzo} \\
\arrayrulecolor{hulinea}\midrule
\endhead
HU-67 & Entrar al desafío & Estudiante & 1 & 5 \\
HU-68 & Leer las reglas antes de empezar & Estudiante & 1 & 3 \\
HU-69 & Responder las preguntas & Estudiante & 1 & 8 \\
HU-70 & No perder las respuestas & Estudiante & 1 & 5 \\
HU-71 & Controlar el tiempo & Estudiante & 1 & 5 \\
HU-72 & Entregar la prueba & Estudiante & 1 & 3 \\
HU-73 & Seguir al grupo durante la prueba & Maestro & 2 & 5 \\
HU-74 & Corregir las pruebas automáticamente & Administrador & 1 & 13 \\
HU-75 & Calcular los puntajes y el ranking & Administrador & 1 & 13 \\
HU-76 & Publicar los resultados & Administrador & 1 & 5 \\
HU-77 & Ver los resultados propios & Estudiante & 1 & 5 \\
HU-78 & Revisar los resultados de un grupo & Maestro & 2 & 8 \\
HU-79 & Revisar los resultados del desafío & Administrador & 1 & 3 \\
HU-80 & Conocer a los mejores de cada categoría & Visitante & 2 & 5 \\
HU-81 & Enviar los resultados para los certificados & Administrador & 2 & 8 \\
HU-82 & Imprimir el cuadernillo de preguntas & Maestro & 1 & 8 \\
HU-83 & Imprimir las hojas de respuestas & Maestro & 1 & 5 \\
HU-84 & Cargar las hojas de respuestas & Maestro & 1 & 8 \\
\arrayrulecolor{hulinea}\midrule
\multicolumn{4}{r}{\textbf{Total de puntos}} & \textbf{115} \\
\arrayrulecolor{hulinea}\bottomrule
\rowcolor{white}\caption{Historias de usuario del módulo de evaluación}\label{tab:historias_evaluacion}\\
\end{longtable}
\arrayrulecolor{black}
}

\subsection{Plan de entregas y alcance del módulo}

\begin{table}[H]
\estilotabla
\begin{tabularx}{\linewidth}{l >{\raggedright\arraybackslash}X >{\raggedright\arraybackslash}p{3.6cm}}
\arrayrulecolor{huazul}\specialrule{1.4pt}{0pt}{0pt}
\rowcolor{hufondo}\textbf{Iteración} & \textbf{Alcance} & \textbf{Historias} \\
\arrayrulecolor{hulinea}\midrule
Iteración 7 (31 de agosto al 11 de septiembre) & Rendición: entrada con código, reglas, respuestas guardadas solas, tiempo, entrega y seguimiento del grupo & \mbox{HU-67} a \mbox{HU-73} \\
Iteración 8 (14 al 25 de septiembre) & Corrección, ranking, resultados del estudiante, del grupo y del desafío, publicación y datos para los certificados & \mbox{HU-74} a \mbox{HU-81} \\
Integración (28 de septiembre al 9 de octubre) & Prueba en papel: cuadernillo, hojas de respuestas y carga de las hojas & \mbox{HU-82} a \mbox{HU-84} \\
\arrayrulecolor{hulinea}\bottomrule
\end{tabularx}
\arrayrulecolor{black}
\caption[Plan de entregas del módulo de evaluación]{Distribución de historias del módulo de evaluación por iteración}
\label{tab:iteraciones_evaluacion}
\end{table}

\section{ITERACIÓN 7}

La Iteración 7 construye la rendición en línea del estudiante. Se presentan su planificación, su diseño, su codificación, sus pruebas y sus resultados.

\subsection{Planificación de la iteración}

La iteración debía garantizar que un estudiante rinda de principio a fin aunque se le corte la conexión, recargue la página o cambie de dispositivo: al volver recupera todo lo que llegó al servidor, y si un guardado falla, lo sabe.

La Tabla~\ref{tab:hus_iteracion_7} lista las historias de la iteración; sus tarjetas, con los criterios de aceptación, están en el Anexo A.

\begin{table}[H]
\estilotabla
\begin{tabularx}{\linewidth}{l >{\raggedright\arraybackslash}X >{\raggedright\arraybackslash}p{2.8cm} c c}
\arrayrulecolor{huazul}\specialrule{1.4pt}{0pt}{0pt}
\rowcolor{hufondo}\textbf{ID} & \textbf{Historia} & \textbf{Usuario} & \textbf{Prioridad} & \textbf{Esfuerzo} \\
\arrayrulecolor{hulinea}\midrule
HU-67 & Entrar al desafío & Estudiante & 1 & 5 \\
HU-68 & Leer las reglas antes de empezar & Estudiante & 1 & 3 \\
HU-69 & Responder las preguntas & Estudiante & 1 & 8 \\
HU-70 & No perder las respuestas & Estudiante & 1 & 5 \\
HU-71 & Controlar el tiempo & Estudiante & 1 & 5 \\
HU-72 & Entregar la prueba & Estudiante & 1 & 3 \\
HU-73 & Seguir al grupo durante la prueba & Maestro & 2 & 5 \\
\arrayrulecolor{hulinea}\midrule
\rowcolor{white}\multicolumn{4}{r}{\textbf{Total de puntos}} & \textbf{34} \\
\arrayrulecolor{hulinea}\bottomrule
\end{tabularx}
\arrayrulecolor{black}
\caption[Historias de la Iteración 7]{Historias de la Iteración 7}
\label{tab:hus_iteracion_7}
\end{table}


\subsection{Diseño}

Cada participante tiene un intento, que nace pendiente al inscribirse. La Figura~\ref{fig:estados_intento} muestra sus tres estados. El intento pasa de pendiente a en curso cuando el estudiante empieza, y de en curso a terminado cuando entrega, se acaba el tiempo o se cierra la rendición. Un intento terminado no vuelve a abrirse.

\begin{figure}[H]
\centering
\begin{tikzpicture}[ancho/.style={minimum width=2.6cm}]
\node[inicio] (ini) at (-3.3,0) {};
\node[estespera, ancho] (pendiente) at (0,0) {Pendiente};
\node[estactivo, ancho] (curso) at (4.8,0) {En curso};
\node[estfinal, ancho, line width=1.6pt] (terminado) at (9.6,0) {Terminado};
\draw[flecha] (ini) -- node[etq, above=2pt]{inscripción} (pendiente);
\draw[flecha] (pendiente) -- node[etq, above=2pt]{empezar} (curso);
\draw[flecha] (curso) -- node[etq, below=3pt]{entrega, fin\\del tiempo\\o cierre} (terminado);
\end{tikzpicture}
\caption{Estados del intento del estudiante}
\label{fig:estados_intento}
\end{figure}

\begin{itemize}
    \item \textbf{Tiempo.} Al empezar, el servidor fija el fin del intento: el menor entre el inicio más la duración y el cierre de la rendición. La cuenta regresiva usa la hora del servidor, no la del dispositivo.
    \item \textbf{Sesión única.} El código personal abre una sesión que se mantiene con una señal cada diez segundos. Entrar desde otro dispositivo cierra la anterior, así el mismo estudiante no rinde dos veces a la vez.
    \item \textbf{Guardado.} Cada respuesta se guarda en cuanto se da y se corrige en ese momento. Si un guardado falla, la pantalla lo avisa y lo reintenta; el estudiante no puede entregar con respuestas sin guardar. Al vencer el tiempo, el intento se cierra con lo que ya está guardado.
    \item \textbf{Privacidad.} El estudiante recibe cada tarea sin su clave ni su explicación. Con la opción de mezclar activada, cada intento ve las opciones en su propio orden, siempre el mismo.
\end{itemize}

\subsection{Codificación}

La pantalla de rendición no muestra la cabecera ni el pie del sitio. Antes de empezar presenta las reglas: la categoría, los minutos, la cantidad de preguntas y cuánto suma o resta cada dificultad, junto con el nombre del estudiante y la opción «¿No eres tú?». Las preguntas se responden de a una o todas juntas, según el desafío, con los mismos componentes del probador y de la práctica. Si el administrador pausa el desafío, la pantalla queda en pausa y al reanudar el tiempo se devuelve.

Mientras tanto, el maestro ve en la pestaña Resultados de su grupo quién está rindiendo y cuántas preguntas lleva, con una actualización cada quince segundos.

\subsection{Pruebas}

Además de los casos de la Tabla~\ref{tab:pruebas_i7}, la rendición se sometió a una prueba de estrés local con 1600 estudiantes simulados que empiezan juntos, descrita en el capítulo 9.

\begin{table}[H]
\estilotabla
\begin{tabularx}{\linewidth}{>{\raggedright\arraybackslash}p{3.6cm} >{\raggedright\arraybackslash}X c c}
\arrayrulecolor{huazul}\specialrule{1.4pt}{0pt}{0pt}
\rowcolor{hufondo}\textbf{Tema} & \textbf{Qué se comprueba} & \textbf{Casos} & \textbf{Resultado} \\
\arrayrulecolor{hulinea}\midrule
Entrada y sesión & Entrada con el código personal, sesión única y hora del servidor antes de empezar & 3 & Aprobado \\
Respuestas & Responder pregunta por pregunta, guardados en orden y entrega bloqueada con respuestas sin guardar & 4 & Aprobado \\
Tiempo y pausa & Cierre automático al acabar el tiempo y vista en pausa & 2 & Aprobado \\
Celular & La rendición ocupa toda la pantalla & 1 & Aprobado \\
\arrayrulecolor{hulinea}\bottomrule
\end{tabularx}
\arrayrulecolor{black}
\caption[Pruebas de la Iteración 7]{Pruebas de extremo a extremo de la Iteración 7}
\label{tab:pruebas_i7}
\end{table}

\subsection{Resultados}

La Figura~\ref{fig:evaluacion_rendicion} muestra la rendición. En (a), la pantalla de reglas que el estudiante lee antes de empezar. En (b), una pregunta durante la prueba, con el tiempo a la vista.

\parcapturas{assets/ui/desafio/rendir-reglas.jpg}{Reglas antes de empezar}{assets/ui/desafio/rendir-pregunta.jpg}{Pregunta durante la rendición}{Rendición del desafío}{fig:evaluacion_rendicion}

\section{ITERACIÓN 8}

La Iteración 8 construye los resultados y la prueba en papel. Se presentan su planificación, su diseño, su codificación, sus pruebas y sus resultados.

\subsection{Planificación de la iteración}

La iteración convirtió las respuestas en resultados: puntajes, ranking por categoría y lo que ve cada usuario. La prueba en papel, para que un aula sin conexión participe con las mismas preguntas y el mismo cálculo, se agregó después, durante la integración; se describe aquí porque usa los mismos intentos, la misma corrección y los mismos resultados.

La Tabla~\ref{tab:hus_iteracion_8} lista las historias de la iteración; sus tarjetas, con los criterios de aceptación, están en el Anexo A.

\begin{table}[H]
\estilotabla
\begin{tabularx}{\linewidth}{l >{\raggedright\arraybackslash}X >{\raggedright\arraybackslash}p{2.8cm} c c}
\arrayrulecolor{huazul}\specialrule{1.4pt}{0pt}{0pt}
\rowcolor{hufondo}\textbf{ID} & \textbf{Historia} & \textbf{Usuario} & \textbf{Prioridad} & \textbf{Esfuerzo} \\
\arrayrulecolor{hulinea}\midrule
HU-74 & Corregir las pruebas automáticamente & Administrador & 1 & 13 \\
HU-75 & Calcular los puntajes y el ranking & Administrador & 1 & 13 \\
HU-76 & Publicar los resultados & Administrador & 1 & 5 \\
HU-77 & Ver los resultados propios & Estudiante & 1 & 5 \\
HU-78 & Revisar los resultados de un grupo & Maestro & 2 & 8 \\
HU-79 & Revisar los resultados del desafío & Administrador & 1 & 3 \\
HU-80 & Conocer a los mejores de cada categoría & Visitante & 2 & 5 \\
HU-81 & Enviar los resultados para los certificados & Administrador & 2 & 8 \\
\arrayrulecolor{hulinea}\midrule
\rowcolor{white}\multicolumn{4}{r}{\textbf{Total de puntos}} & \textbf{60} \\
\arrayrulecolor{hulinea}\bottomrule
\end{tabularx}
\arrayrulecolor{black}
\caption[Historias de la Iteración 8]{Historias de la Iteración 8}
\label{tab:hus_iteracion_8}
\end{table}


\subsection{Diseño}

La Figura~\ref{fig:flujo_resultados} resume el camino de una respuesta hasta los resultados, de izquierda a derecha. Cada respuesta se corrige al guardarla; al cerrar la rendición se calculan los puntajes y el ranking de cada categoría; los resultados se publican al llegar la fecha y los ven el estudiante, el maestro, la portada y los certificados.

\begin{figure}[H]
\centering
\begin{tikzpicture}[ancho/.style={minimum width=2.9cm, text width=2.6cm}]
\node[servicio, ancho] (guardar) at (0,0) {Cada respuesta se corrige al guardarla};
\node[servicio, ancho] (cerrar) at (4.1,0) {Cierre de la rendición};
\node[datos, minimum width=2.9cm] (calcular) at (8.2,0) {Puntajes y\\ranking por\\categoría};
\node[servicio, ancho] (publicar) at (12.3,0) {Publicación al llegar la fecha};
\node[interfaz, minimum width=3.6cm] (portada) at (12.3,-2.6) {Estudiante, maestro,\\portada y certificados};
\draw[flecha] (guardar) -- (cerrar);
\draw[flecha] (cerrar) -- (calcular);
\draw[flecha] (calcular) -- (publicar);
\draw[flecha] (publicar) -- (portada);
\end{tikzpicture}
\caption{De las respuestas a los resultados}
\label{fig:flujo_resultados}
\end{figure}

\begin{itemize}
    \item \textbf{Puntaje.} El puntaje de un estudiante es el puntaje inicial de su categoría más lo que suma o resta cada pregunta según su dificultad (Tabla~\ref{tab:puntaje}); una pregunta sin responder no resta.
    \item \textbf{Ranking.} Se ordena por categoría: primero el mayor puntaje y, si hay empate, el menor tiempo usado.
    \item \textbf{Publicación.} Los resultados se publican solos una sola vez: desde la fecha de resultados (o, sin fecha, desde el cierre de la rendición), la primera consulta que llega al servidor los calcula y los publica; el administrador puede ocultarlos y volver a publicarlos. Quedan visibles en la portada hasta el fin del periodo de resultados.
    \item \textbf{Prueba en papel.} El maestro imprime un cuadernillo por categoría y una hoja de respuestas por estudiante, con su nombre y su código personal. Todo se marca en la hoja: el cuadernillo no lleva respuestas, así que se fotocopia y se vuelve a usar. Cada tipo de pregunta se adapta al papel con letras, números o una marca sobre la imagen (Tabla~\ref{tab:papel}), y la pantalla de carga usa las mismas letras y números, para que lo marcado se traslade sin dudas.
    \item \textbf{Intento en papel.} Lo cargado se guarda como un intento más, marcado como rendido en papel, y se corrige igual que uno en línea. Cada estudiante rinde de una sola forma: con su hoja cargada no puede entrar en línea, y al revés. Como el papel no mide el tiempo, en un empate el intento en línea va primero.
    \item \textbf{Plazos.} La prueba se imprime desde dos días antes de su inicio, para dar tiempo a fotocopiarla sin adelantar demasiado las preguntas. Las hojas se cargan desde el inicio hasta que se calculan los resultados.
    \item \textbf{Qué ve cada uno.} El administrador decide qué ve el estudiante al entregar y qué al publicar: su puntaje, cuáles acertó y las soluciones. Su puesto solo aparece con los resultados publicados. En lo público, un menor aparece con su nombre y la inicial del apellido, su colegio y su ciudad.
\end{itemize}

\begin{table}[H]
\estilotabla
\begin{tabularx}{\linewidth}{>{\raggedright\arraybackslash}p{3.2cm} >{\raggedright\arraybackslash}X >{\raggedright\arraybackslash}X}
\arrayrulecolor{huazul}\specialrule{1.4pt}{0pt}{0pt}
\rowcolor{hufondo}\textbf{Tipo de pregunta} & \textbf{En el cuadernillo} & \textbf{En la hoja de respuestas} \\
\arrayrulecolor{hulinea}\midrule
Opción múltiple & Las opciones con su letra & Un círculo por letra, para rellenar \\
Respuesta corta & El enunciado & Un recuadro para escribir \\
Arrastrar y soltar & La imagen con cada lugar numerado y las piezas con su letra & Una casilla por lugar, para la letra de la pieza \\
Zonas sobre imagen & La imagen completa & Una copia pequeña de la imagen, para marcar una X \\
Estados por casilla & Los estados con su letra & La cuadrícula, para escribir una letra en cada casilla \\
Huecos en el texto & El texto con cada hueco numerado y las opciones con su letra & Una casilla por hueco \\
\arrayrulecolor{hulinea}\bottomrule
\end{tabularx}
\arrayrulecolor{black}
\caption[Adaptación de las preguntas al papel]{Cómo se presenta y se responde cada tipo de pregunta en papel}
\label{tab:papel}
\end{table}

\subsection{Codificación}

Como cada respuesta ya está corregida, cerrar todos los intentos abiertos y calcular los puntajes de un desafío completo son dos operaciones sobre la base de datos, y el ranking de todas las categorías, una sola; la cantidad de consultas no crece con la cantidad de estudiantes. Esta forma de cálculo reemplazó a una primera versión que recalculaba el ranking en cada entrega, y que con unos mil estudiantes habría superado el límite de consultas de Cloudflare D1.

El maestro ve los resultados de su grupo con el puntaje, las correctas, el puesto y cada respuesta, la pregunta que más costó, y los descarga en Excel. El administrador ve los de todo el desafío por categoría. La portada muestra los diez primeros de cada categoría, y una exportación protegida con clave entrega los resultados al CMS para los certificados (capítulo 9).

La prueba en papel se arma en el navegador con las mismas tareas que se rinden en línea, sin sus claves. Cada pregunta ocupa una hoja tamaño carta: si no cabe, el sistema reduce poco a poco sus imágenes, y solo las preguntas con mucho texto siguen en una segunda hoja, cortadas entre párrafos. La portada muestra el castor del departamento del colegio del maestro, las reglas, el puntaje y ejemplos de cómo marcar, solo de los tipos de pregunta que trae la prueba. El maestro la imprime o la descarga en PDF para llevarla a una fotocopiadora. Para cargar las hojas, una pantalla recorre a los estudiantes del grupo en orden, con controles que reproducen la hoja; «Guardar y seguir» pasa al siguiente que tenga su hoja pendiente.

\subsection{Pruebas}

\begin{table}[H]
\estilotabla
\begin{tabularx}{\linewidth}{>{\raggedright\arraybackslash}p{3.6cm} >{\raggedright\arraybackslash}X c c}
\arrayrulecolor{huazul}\specialrule{1.4pt}{0pt}{0pt}
\rowcolor{hufondo}\textbf{Tema} & \textbf{Qué se comprueba} & \textbf{Casos} & \textbf{Resultado} \\
\arrayrulecolor{hulinea}\midrule
Publicación & Publicación automática al cierre, también de quien no entregó; ocultar y volver a publicar; periodo de resultados & 4 & Aprobado \\
Ranking & Desempate por tiempo y ranking público con nombre e inicial & 2 & Aprobado \\
Resultados & Resultados del grupo y del desafío, con sus descargas & 2 & Aprobado \\
Prueba en papel & Impresión desde dos días antes, una hoja por estudiante y PDF tamaño carta; carga, corrección y conteo en los resultados; exclusión con la rendición en línea; hoja en blanco, cambios sin guardar y quitar lo cargado; plazos y acceso de otro maestro & 4 & Aprobado \\
\arrayrulecolor{hulinea}\bottomrule
\end{tabularx}
\arrayrulecolor{black}
\caption[Pruebas de la Iteración 8]{Pruebas de extremo a extremo de la Iteración 8}
\label{tab:pruebas_i8}
\end{table}

La prueba de estrés del capítulo 9 verificó además, al cierre, los 1600 puntajes frente a las respuestas enviadas y que los puestos de cada categoría fueran consecutivos y ordenados por puntaje; el desempate por tiempo tiene su propia prueba.

Además, se comprobó en los cuadernillos de las seis categorías que cada pregunta empieza en su propia hoja, que solo las de mucho texto siguen en la siguiente, y que la cantidad de hojas es la misma en pantalla, en el PDF y al imprimir.

\subsection{Resultados}

La Figura~\ref{fig:evaluacion_resultados_grupo} muestra lo que se ve al terminar. En (a), la pantalla final del estudiante, con lo que el desafío permite mostrarle. En (b), los resultados del grupo, que el maestro revisa en la pestaña Resultados.

\parcapturas{assets/ui/desafio/rendir-final.jpg}{Pantalla final del estudiante}{assets/ui/desafio/grupo-resultados.jpg}{Resultados del grupo}{Resultados del estudiante y del grupo}{fig:evaluacion_resultados_grupo}

En la Figura~\ref{fig:evaluacion_resultados} aparecen los resultados del desafío. En (a), el administrador los revisa por categoría. En (b), el público consulta el ranking en la portada.

\parcapturas{assets/ui/desafio/desafio-resultados.jpg}{Resultados por categoría}{assets/ui/desafio/portada-ranking.jpg}{Ranking público}{Resultados del desafío}{fig:evaluacion_resultados}

Para las aulas sin conexión, el maestro imprime la prueba de su grupo. La Figura~\ref{fig:evaluacion_papel_hojas} muestra las tres piezas: la portada del cuadernillo, una pregunta por hoja y la hoja de respuestas de cada estudiante.

\begin{figure}[htbp]
\centering
\captura{assets/ui/desafio/papel-hojas.jpg}
\caption{Portada del cuadernillo, una pregunta y una hoja de respuestas}
\label{fig:evaluacion_papel_hojas}
\end{figure}

La Figura~\ref{fig:evaluacion_papel} muestra el resto del proceso. En (a) está la pantalla para imprimir la prueba. Después de la prueba, el maestro copia en la pantalla de carga, en (b), lo que marcó cada estudiante.

\parcapturas{assets/ui/desafio/papel-imprimir.jpg}{Impresión de la prueba}{assets/ui/desafio/papel-carga.jpg}{Carga de una hoja de respuestas}{Prueba en papel}{fig:evaluacion_papel}
